% =====================================================================
%  setup/problems.tex -- solved problems, practice lists, difficulty tags
%  and automatic problem statistics.
%
%  SOLVED PROBLEM (header; the solution follows as normal text):
%    \solvedproblem{c01-max-depth}{Maximum Depth}{LeetCode 104}
%        {https://leetcode.com/problems/maximum-depth-of-binary-tree/}{LC Easy}
%        {recursion, DFS on a binary tree}
%    -> "W1.1.1  Maximum Depth", label prob:c01-max-depth, \probref{c01-max-depth}
%    Inside the solution use \solstep{...} headings and the observation env.
%
%  PRACTICE LIST (after a topic):          CHAPTER-END MIXED LIST:
%    \begin{practiceset}                      \begin{mixedset}
%      \practice{Title}{Source}{URL}{LC Easy}{hint}   (same \practice macro)
%    \end{practiceset}                        \end{mixedset}
%    -> Q1.3.1, Q1.3.2, ...                   -> M1.1, M1.2, ...
%
% Difficulty strings: "CF 1300" (coloured by rating), "LC Easy|Medium|Hard",
% anything else (e.g. "CSES", "AtCoder 1580", "Kattis") is shown in grey.
% LEVEL BADGE: end the string with a plain " L1".." L6" suffix (NOT \lv{} --
% xstring expands macros and chokes); it is peeled off and drawn as a second
% chip in the ladder colour lv1..lv6.  URLs must not contain  %  #  &.
% =====================================================================
\makeatletter

% ---------------- Statistics (recorded to the .aux, shown via \stat) ----------------
\newcounter{chproblems}[chapter]
\newcounter{chsolved}[chapter]
\newcounter{partproblems}[part]
\newcounter{partchapters}[part]
\newcounter{bookproblems}
\newcounter{booksolved}
\newcounter{bookwritten}
\newcounter{bookplanned}
\newcounter{writtenchapters}

\newcommand{\countwritten}{%
  \stepcounter{chproblems}\stepcounter{partproblems}%
  \stepcounter{bookproblems}\stepcounter{bookwritten}}
\newcommand{\countplanned}[1]{%
  \addtocounter{chproblems}{#1}\addtocounter{partproblems}{#1}%
  \addtocounter{bookproblems}{#1}\addtocounter{bookplanned}{#1}}

% \recordstat{key}{value}  stores a number;  \stat{key}  prints it (next run).
% \recordstat writes when the page is shipped out (like \label). At the end of the
% document the last page is already gone, so \recordstatnow writes immediately.
\newcommand{\recordstat}[2]{\begingroup\edef\@currentlabel{#2}\label{stat:#1}\endgroup}
\newcommand{\recordstatnow}[2]{%
  \immediate\write\@auxout{\string\newlabel{stat:#1}{{#2}{}{}{}{}}}}
\newcommand{\stat}[1]{\ifcsname r@stat:#1\endcsname\getrefnumber{stat:#1}\else\textbf{?}\fi}

\AtEndDocument{%
  \recordstatnow{book:problems}{\arabic{bookproblems}}%
  \recordstatnow{book:solved}{\arabic{booksolved}}%
  \recordstatnow{book:written}{\arabic{bookwritten}}%
  \recordstatnow{book:planned}{\arabic{bookplanned}}%
  \recordstatnow{book:writtenchapters}{\arabic{writtenchapters}}}

% ---------------- Difficulty tag ----------------
\newcommand{\dt@cfcolor}[1]{%
  \ifnum#1<1200 \def\dt@col{diffgray}\else
  \ifnum#1<1400 \def\dt@col{diffgreen}\else
  \ifnum#1<1600 \def\dt@col{diffcyan}\else
  \ifnum#1<1900 \def\dt@col{diffblue}\else
  \ifnum#1<2100 \def\dt@col{diffviolet}\else
  \ifnum#1<2400 \def\dt@col{difforange}\else
  \def\dt@col{diffred}\fi\fi\fi\fi\fi\fi}
\newcommand{\difficultytag}[1]{%
  \begingroup
  % ---- peel a trailing " L<digit>" level suffix BEFORE any xstring parsing ----
  \def\dt@lv{0}\def\dt@main{#1}%
  \StrLen{#1}[\dt@len]%
  \ifnum\dt@len>2\relax
    \StrRight{#1}{3}[\dt@tail]%
    \IfBeginWith{\dt@tail}{ L}{%
      \StrGobbleRight{#1}{3}[\dt@main]%
      \StrChar{#1}{\dt@len}[\dt@lv]%
    }{}%
  \fi
  % ---- colour the main chip from the (now macro-free) remaining string ----
  \def\dt@col{darkgray}%
  \IfBeginWith{\dt@main}{CF }{\StrBehind{\dt@main}{CF }[\dt@num]\IfInteger{\dt@num}{\dt@cfcolor{\dt@num}}{}}{}%
  \IfBeginWith{\dt@main}{LC }{%
    \IfSubStr{\dt@main}{Easy}{\def\dt@col{diffgreen}}{}%
    \IfSubStr{\dt@main}{Medium}{\def\dt@col{difforange}}{}%
    \IfSubStr{\dt@main}{Hard}{\def\dt@col{diffred}}{}}{}%
  \tikz[baseline=(dt.base)]\node[fill=\dt@col, text=white, rounded corners=1.5pt,
    inner xsep=3pt, inner ysep=1.6pt, font=\scriptsize\bfseries\sffamily](dt){\dt@main};%
  \ifnum\dt@lv=0\else
    \hspace{2pt}\tikz[baseline=(dl.base)]\node[fill=lv\dt@lv, text=white,
      rounded corners=1.5pt, inner xsep=2.5pt, inner ysep=1.6pt,
      font=\scriptsize\bfseries\sffamily](dl){L\dt@lv};%
  \fi
  \endgroup}

% ---------------- Solved problems ----------------
\newcounter{solvedprob}[section]
\renewcommand{\thesolvedprob}{\thesection.\arabic{solvedprob}}
\newcounter{observation}[solvedprob]
\newcommand{\currentchaptertitle}{}

% Header drawing only (no counting); used by \solvedproblem and by the legend.
\newcommand{\solvedheader}[6]{%  #1 id  #2 title  #3 source  #4 url  #5 difficulty  #6 techniques
  \begin{tcolorbox}[enhanced, colback=primaryblue!5, colframe=primaryblue, coltext=textonlight,
      boxrule=1.3pt, arc=3mm, left=8pt, right=8pt, top=3pt, bottom=3pt,
      title={#1\quad #2}, fonttitle=\bfseries\large,
      coltitle=textondark, colbacktitle=primaryblue]
    \small\textbf{Source:}~\href{#4}{#3}\hfill\difficultytag{#5}\par
    \textbf{Techniques:}~#6
  \end{tcolorbox}}

\newcommand{\solvedproblem}[6]{%  #1 label  #2 title  #3 source  #4 url  #5 difficulty  #6 techniques
  \par\needspace{12\baselineskip}\medskip
  \ifnum\value{chsolved}=0
    \addtocontents{lop}{\protect\lopchapter{\thechapter}{\currentchaptertitle}}%
  \fi
  \refstepcounter{solvedprob}\label{prob:#1}%
  \countwritten\stepcounter{chsolved}\stepcounter{booksolved}%
  \addcontentsline{lop}{problem}{\protect\numberline{W\thesolvedprob}#2~--~#3}%
  \solvedheader{W\thesolvedprob}{#2}{#3}{#4}{#5}{#6}\par\nobreak}

% Run-in step heading inside a solution:  \solstep{Key Observations}
\newcommand{\solstep}[1]{%
  \par\needspace{4\baselineskip}\medskip
  \noindent{\large\bfseries\color{deepblue}\ding{228}\ #1}\par\nobreak\smallskip}

% Numbered observations inside a solved problem (restart at every problem).
\newenvironment{observation}
  {\par\refstepcounter{observation}\noindent
   \textbf{\textcolor{accentorange}{Observation \arabic{observation}.}}\ \ignorespaces}
  {\par}

% ---------------- Practice and mixed lists ----------------
\newcounter{practiceprob}[section]
\renewcommand{\thepracticeprob}{\thesection.\arabic{practiceprob}}
\newcounter{mixedprob}[chapter]
\renewcommand{\themixedprob}{\thechapter.\arabic{mixedprob}}
\newcommand{\practice@prefix}{Q}
\newcommand{\practice@counter}{practiceprob}
\newcommand{\practice@listsetup}{%
  \setlength{\leftmargin}{4.6em}\setlength{\labelwidth}{4.3em}%
  \setlength{\labelsep}{0.3em}\setlength{\itemsep}{0.6em}%
  \setlength{\parsep}{0.15em}\setlength{\topsep}{0.5em}}

\newenvironment{practiceset}
  {\renewcommand{\practice@prefix}{Q}\renewcommand{\practice@counter}{practiceprob}%
   \begin{list}{}{\practice@listsetup}}
  {\end{list}}
\newenvironment{mixedset}
  {\renewcommand{\practice@prefix}{M}\renewcommand{\practice@counter}{mixedprob}%
   \begin{list}{}{\practice@listsetup}}
  {\end{list}}

% Entry drawing only (no counting); used by \practice and by the legend.
\newcommand{\practiceitem}[6]{%  #1 id  #2 title  #3 source  #4 url  #5 difficulty  #6 hint
  \item[\textbf{\textcolor{primaryblue}{#1}}]%
  \textbf{#2}\ \textcolor{darkgray}{--}\ \href{#4}{#3}\hfill\difficultytag{#5}\par
  {\small\color{darkgray}\textit{Hint:}\ #6\par}}

\newcommand{\practice}[5]{%  #1 title  #2 source  #3 url  #4 difficulty  #5 hint
  \refstepcounter{\practice@counter}\countwritten
  \practiceitem{\practice@prefix\csname the\practice@counter\endcsname}{#1}{#2}{#3}{#4}{#5}}

% Progressive hint ladder: a gentle nudge, then a stronger one (optional tool).
% Use inside practiceset/mixedset exactly like \practice, but with two hints.
\newcommand{\practicehh}[6]{%  #1 title #2 source #3 url #4 difficulty #5 hint1 #6 hint2
  \refstepcounter{\practice@counter}\countwritten
  \item[\textbf{\textcolor{primaryblue}{\practice@prefix\csname the\practice@counter\endcsname}}]%
  \textbf{#1}\ \textcolor{darkgray}{--}\ \href{#3}{#2}\hfill\difficultytag{#4}\par
  {\small\color{darkgray}\textit{Hint 1:}\ #5\par
   \textit{Hint 2:}\ #6\par}}

% ---------------- List of solved problems ----------------
\newcommand{\l@problem}{\@dottedtocline{1}{0em}{4.6em}}
\newcommand{\lopchapter}[2]{%
  \par\addvspace{0.9em}\noindent{\bfseries\color{primaryblue}Chapter #1\quad #2}\par\nopagebreak\smallskip}
\newcommand{\listofsolvedproblems}{\@starttoc{lop}}

\makeatother
